% !TEX program = xelatex
% !BIB program = biber
\documentclass{beamer}
\usepackage[english]{babel}
\usepackage[orientation=portrait,size=a0,scale=1.4]{beamerposter}
\usepackage[IWES]{tensygrid}
\usepackage[backend=biber]{biblatex}
\usepackage{lipsum}
\author{Luke Skywalker}
\title{Fraunhofer {\LaTeX} Poster Template}
\subtitle{Combine the Beauty of Science with Excellent Typesetting}
\date[International Symposium on Typography 2020]{February 17, 2020}
\begin{document}
\begin{frame}
\maketitle
\begin{columns}[onlytextwidth,t]
\begin{column}{.30\textwidth}
  \begin{block}{Introduction}
    The basic \structure{block} looks like this.
  \end{block}
  \begin{exampleblock}{Does it come in black?}
    Sure, use an \textbf{exampleblock}!
  \end{exampleblock}
  \begin{alertblock}{How do you make it pop?}
    Use an \alert{alertblock}!
  \end{alertblock}
  \begin{block}{Methods}
    \lipsum[1]
  \end{block}
\end{column}
\begin{column}{.30\textwidth}
  \begin{block}{Results}
    \lipsum[2]
    \unskip
  \end{block}
  \begin{block}{Conclusions}
    \lipsum[4]
  \end{block}
\end{column}
\begin{column}{.30\textwidth}
  \begin{block}{Acknowledgements}
    \lipsum[5]
  \end{block}
  \begin{block}{References}
    \lipsum[6]
  \end{block}
  \begin{block}{Contact information}
    \lipsum[7]
  \end{block}
\end{column}
\end{columns}
\end{frame}
\end{document}
